\documentclass[12pt,a4paper]{article}
\usepackage[margin=3cm]{geometry}
\usepackage{newtxtext,newtxmath} 

\begin{document}

% ====================
% HALAMAN SAMPUL
% ====================
\thispagestyle{empty}

\begin{center}
  \textbf{\LARGE KAJIAN PUSTAKA TOPIK TESIS}\\[0.5cm]
  \textit{Mata Kuliah Metode Penelitian}
\end{center}

\vfill

\noindent
\textbf{Disusun oleh:}\\[0.3cm]
\begin{tabular}{@{}p{1.2cm}p{7.5cm}@{}}
  Nama  & : \dotfill \\
  NRP   & : \dotfill \\
  Kelas & : \dotfill \\
\end{tabular}

\bigskip\bigskip\bigskip

\noindent
\textbf{Dosen Pengampu:}\\[0.3cm]
\mbox{}\dotfill

\vfill

\begin{center}
  \textbf{Program Studi S2 Matematika}\\
  \textbf{Departemen Matematika}\\
  \textbf{Fakultas Sains dan Analitika Data}\\
  \textbf{Institut Teknologi Sepuluh Nopember}\\
  \textbf{2026}
\end{center}

\newpage
\setcounter{page}{1}

% ====================
% BAB 1
% ====================
\noindent\textbf{\Large BAB 1 PENDAHULUAN}\par\bigskip

\noindent\textbf{\large 1.1 Latar Belakang}\par
Tuliskan latar belakang singkat mengenai topik yang dipilih. Jelaskan mengapa topik tersebut penting, menarik, atau relevan untuk dikaji lebih lanjut dalam tesis.\par\bigskip

\noindent\textbf{\large 1.2 Topik Tesis}\par
Tuliskan kandidat topik/judul tesis yang akan anda kerjakan. Jelaskan pula \textbf{kebaruan} dan \textbf{kebermanfaatan} dari topik yang diusulkan.

\newpage
% ====================
% BAB 2
% ====================
\noindent\textbf{\Large BAB 2 KAJIAN ARTIKEL}\par\bigskip

\noindent\textbf{\large 2.1 Artikel 1}\par\bigskip

\noindent Identitas Artikel
\begin{itemize}
  \item[-] Penulis: \dotfill
  \item[-] Tahun: \dotfill
  \item[-] Judul: \dotfill
  \item[-] Jurnal/Prosiding: \dotfill
  \item[-] DOI/Link (jika ada): \dotfill
\end{itemize}

\bigskip
\noindent Ringkasan Singkat Artikel\par
Tuliskan ringkasan singkat isi artikel ini dalam 1 paragraf.\par
\noindent\mbox{}\dotfill\par
\noindent\mbox{}\dotfill\par
\noindent\mbox{}\dotfill

\bigskip
\noindent Analisis Artikel
\begin{enumerate}
  \item \textbf{Kaitan dengan topik yang dipilih}\par
        \noindent\mbox{}\dotfill\par
        \noindent\mbox{}\dotfill
  \item \textbf{Kontribusi artikel}\par
        \noindent\mbox{}\dotfill\par
        \noindent\mbox{}\dotfill
  \item \textbf{Novelty artikel}\par
        \noindent\mbox{}\dotfill\par
        \noindent\mbox{}\dotfill
  \item \textbf{Kekurangan/keterbatasan artikel}\par
        \noindent\mbox{}\dotfill\par
        \noindent\mbox{}\dotfill
\end{enumerate}

\newpage
% Lanjutan BAB 2 (Artikel 2)
\noindent\textbf{\large 2.2 Artikel 2}\par\bigskip

\noindent Identitas Artikel
\begin{itemize}
  \item[-] Penulis: \dotfill
  \item[-] Tahun: \dotfill
  \item[-] Judul: \dotfill
  \item[-] Jurnal/Prosiding: \dotfill
  \item[-] DOI/Link (jika ada): \dotfill
\end{itemize}

\bigskip
\noindent Ringkasan Singkat Artikel\par
Tuliskan ringkasan singkat isi artikel ini dalam 1 paragraf.\par
\noindent\mbox{}\dotfill\par
\noindent\mbox{}\dotfill\par
\noindent\mbox{}\dotfill

\bigskip
\noindent Analisis Artikel
\begin{enumerate}
  \item \textbf{Kaitan dengan topik yang dipilih}\par
        \noindent\mbox{}\dotfill\par
        \noindent\mbox{}\dotfill
  \item \textbf{Kontribusi artikel}\par
        \noindent\mbox{}\dotfill\par
        \noindent\mbox{}\dotfill
  \item \textbf{Novelty artikel}\par
        \noindent\mbox{}\dotfill\par
        \noindent\mbox{}\dotfill
  \item \textbf{Kekurangan/keterbatasan artikel}\par
        \noindent\mbox{}\dotfill\par
        \noindent\mbox{}\dotfill
\end{enumerate}

\newpage
% Lanjutan BAB 2 (Artikel 3 - 8)
\noindent\textbf{\large 2.3 Artikel 3}\par\bigskip

\noindent\textbf{\large 2.4 Artikel 4}\par\bigskip

\noindent\textbf{\large 2.5 Artikel 5}\par\bigskip

\noindent\textbf{\large 2.6 Artikel 6}\par\bigskip

\noindent\textbf{\large 2.7 Artikel 7}\par\bigskip

\noindent\textbf{\large 2.8 Artikel 8}

\newpage
% ====================
% BAB 3
% ====================
\noindent\textbf{\Large BAB 3 SINTESIS KAJIAN PUSTAKA}\par\bigskip

\noindent\textbf{\large 3.1 Perbandingan Antarartikel}\par
Pada bagian ini, bandingkan artikel-artikel yang telah dikaji. Jelaskan persamaan, perbedaan, metode yang digunakan, hasil yang diperoleh, dan arah perkembangan penelitian.\par\bigskip

\noindent\textbf{\large 3.2 Research Gap}\par
Berdasarkan hasil kajian lima artikel, jelaskan celah penelitian yang masih terbuka.\par\bigskip

\noindent\textbf{\large 3.3 State-of-the-Art}\par
Jelaskan perkembangan mutakhir penelitian pada topik yang dipilih. Uraikan pendekatan, metode, atau fokus penelitian yang saat ini paling dominan atau paling maju.

\newpage
% ====================
% BAB 4
% ====================
\noindent\textbf{\Large BAB 4 KESIMPULAN}\par\bigskip

\noindent Berdasarkan kajian pustaka yang telah dilakukan, berikan beberapa hal yang dapat disimpulkan. Kemudian sebutkan rumusan masalah (\textit{research question}) yang cocok dengan topik yang diusulkan serta arah awal penelitian yang mungkin dilakukan.

\newpage
% ====================
% DAFTAR PUSTAKA
% ====================
\noindent\textbf{\Large DAFTAR PUSTAKA}\par\bigskip

\noindent Tuliskan semua referensi yang digunakan dengan format sitasi Harvard style.

\end{document}